\scanpage{66}
% archive foliation: 65
4. Détermination de $M_{11}(4)$.

Soit $E$ une courbe elliptique sur $S$ (à car. résiduelles $\ne2$), considérons
\[
\tag{38}M={}_4E\qquad\begin{array}{l}
\text{schéma en groupes sur }S\\
\text{loc. isom. à }(\mathbb Z/4\mathbb Z)^2,\\
\text{i.e. }\deleted{\text{syst. local}}\text{ faisceau abélien}\\
\text{sur }S_{\text{ét}}\text{ qui est un }(\mathbb Z/4\mathbb Z)_S\\
\text{module loc. libre de rg }2.
\end{array}
\]
de sorte qu'on a
\deleted{(39)}
\[
\tag{\unsure{39|34}}\diagram{\begin{tikzcd}[ampersand replacement=\&]
0\arrow[r]\&M_0\arrow[r]\&M\arrow[r,"\mathrm{can}"]\&M_0\arrow[r]\&0\\
\&M\arrow[u,"\mathrm{can}"]\arrow[ur,"2\operatorname{id}_M"']\&\&\&
\end{tikzcd}}
\]
où
\[
M_0=M\otimes_{\mathbb Z/4\mathbb Z}\mathbb F_2\simeq{}_2E.
\]
La donnée d'un $M$ équivaut à celle d'un $M_0$ (i.e. d'un $\underline J$ \deleted{\unsure{rev.}} étale de degré $3$ de $S$ \deleted{$\unsure{\underline J}$}
\[
\tag{40}\underline J\simeq M_0^*,\qquad M_0=V_{\underline J}(\mathbb F_2)\simeq\mathbb F_2^{\underline J}/\underset{\mathrm{diag}}{\underline{\mathbb F_2}}),
\]
\uline{et} d'un torseur \inserted{$P$} sous
\[
\tag{41}\unsure{N|V}_{\underline J}(M_0)\simeq\underline{\operatorname{End}}_{\mathbb F_2,S}(M_0)\simeq M_0\oplus\mathbb F_2^{\underline\omega}
\]
\margin{$\underline\omega$ est le rev. étale de degré $2$ ``des orientations'' associé au rev. $\underline J$ de degré $3$ de $S$}
Un ``épinglage'' de $M={}_4E$ (i.e. \deleted{d'une} la donnée d'un iso
\[
\tag{42}(\mathbb Z/4\mathbb Z)^2_S\simeq M\;(={}_4E)
\]
équivaut à celle d'une \inserted{trivialisation de $\underline J$, plus d'une} section \unsure{(trivialisé)}-$\unsure{\underline P|\underline E}$
